\chapter{Concluzii}
\pagestyle{fancy}

Lucrarea de față a avut ca obiectiv proiectarea, implementarea și validarea unei platforme pentru monitorizarea consumului de utilități printr-un flux IoT bazat pe LoRaWAN. Sistemul dezvoltat demonstrează integrarea mai multor componente software specifice aplicațiilor distribuite moderne: simulare de dispozitive, server de rețea LoRaWAN, broker MQTT, procesare de telemetrie, stocare relațională și time-series, interfață web operațională și modul de predicție a consumului.

Din punct de vedere funcțional, aplicația permite monitorizarea mai multor categorii de contoare inteligente, corespunzătoare utilităților de tip electricitate, gaz, apă, încălzire și răcire. Datele generate de simulator sunt transmise prin infrastructura ChirpStack, publicate prin MQTT, procesate de worker-ul aplicației și salvate în InfluxDB. Metadatele dispozitivelor, utilizatorii, tarifele și relațiile de ownership sunt gestionate separat în PostgreSQL. Această separare reflectă o decizie arhitecturală importantă, deoarece telemetria are caracter temporal și volum ridicat, în timp ce datele aplicației au caracter relațional.

\section{Contribuții realizate}

Prima contribuție importantă constă în realizarea unui flux complet de date, de la dispozitivul LoRaWAN simulat până la afișarea informațiilor în dashboard. Chiar dacă mediul folosește dispozitive simulate, traseul tehnic reproduce etapele principale ale unei integrări IoT reale: generarea pachetelor, transmiterea prin gateway bridge, gestionarea dispozitivelor în ChirpStack, publicarea mesajelor MQTT și ingestia în baza de date time-series.

A doua contribuție este reprezentată de implementarea aplicației web, care oferă utilizatorului o interfață centralizată pentru vizualizarea flotei de dispozitive. Dashboard-ul include pagini pentru prezentare generală, administrarea contoarelor, analiza unui contor individual, estimarea costurilor și vizualizarea geografică. De asemenea, aplicația diferențiază între conturi de companie și utilizatori individuali, iar asocierea dispozitivelor este realizată printr-un mecanism de revendicare pe bază de cod.

A treia contribuție constă în integrarea unui modul predictiv bazat pe modelul ARIMA. Acesta utilizează datele istorice agregate din InfluxDB pentru a genera estimări pe termen scurt ale consumului. Integrarea acestui modul transformă aplicația dintr-un sistem strict descriptiv într-un sistem care include și analiză predictivă, oferind utilizatorului o perspectivă asupra evoluției probabile a consumului.

O altă contribuție importantă este reprezentată de generarea unui set de date demonstrative și istorice. Acesta permite testarea dashboard-ului, a calculelor de cost și a modulului ARIMA fără dependență de hardware fizic. În contextul unei lucrări de licență, această abordare este utilă deoarece permite reproducerea experimentelor și demonstrarea funcționalității sistemului într-un mediu local controlat.

\section{Analiză critică a rezultatelor}

Rezultatele obținute arată că sistemul implementat îndeplinește obiectivele principale propuse. Fluxul de date este funcțional, dispozitivele pot fi gestionate și asociate utilizatorilor, valorile de consum sunt afișate în interfață, iar predicția ARIMA poate fi calculată pentru seriile istorice disponibile. Testele statice, unitare, de integrare și validările vizuale au confirmat funcționarea componentelor principale.

Cu toate acestea, sistemul are și limitări. Cea mai importantă limitare este folosirea unui simulator în locul unor contoare și gateway-uri fizice. Această alegere permite dezvoltarea și validarea logicii software, dar nu acoperă aspecte precum interferențele radio, pierderile de pachete în condiții reale, autonomia energetică a dispozitivelor sau comportamentul hardware. În plus, tarifele utilizate pentru calculul costurilor sunt statice și au rol demonstrativ, deci rezultatele nu trebuie interpretate ca facturare reală.

O altă limitare este faptul că predicția ARIMA este calculată la cerere și depinde de calitatea datelor istorice disponibile. Modelul este potrivit pentru demonstrarea predicției pe termen scurt, însă performanța sa poate varia în funcție de regularitatea consumului și de comportamentul fiecărui dispozitiv. Pentru un sistem de producție ar fi necesară o evaluare mai amplă a acurateței predicțiilor, folosind date reale și indicatori cantitativi precum MAE, RMSE sau MAPE.

\section{Dezvoltări ulterioare}

O direcție evidentă de dezvoltare este înlocuirea simulatorului cu dispozitive LoRaWAN reale. Acest pas ar permite validarea sistemului în condiții operaționale și ar oferi o evaluare mai realistă a performanței fluxului de comunicație. În același timp, ar putea fi analizate aspecte precum configurarea gateway-urilor, acoperirea radio și comportamentul dispozitivelor în timp.

O altă direcție este extinderea modulului de predicție. Pe lângă ARIMA, sistemul ar putea integra metode precum Holt-Winters, Random Forest sau rețele neuronale recurente, iar rezultatele acestora ar putea fi comparate pe aceleași seturi de date. De asemenea, predicțiile ar putea fi salvate periodic, nu doar calculate la cerere, pentru a permite analiză istorică asupra calității forecast-ului.

Aplicația poate fi extinsă și din perspectiva funcționalităților oferite utilizatorului. Exemple relevante sunt definirea unor alerte pentru consum anormal, exportul rapoartelor, configurarea unor tarife diferențiate pe intervale orare și adăugarea unor roluri administrative mai detaliate. Pentru un scenariu comercial, ar fi necesare și mecanisme suplimentare de securitate, audit și administrare a organizațiilor.

În concluzie, proiectul demonstrează că o platformă locală bazată pe LoRaWAN, MQTT, baze de date specializate și dashboard web poate fi utilizată pentru monitorizarea și analiza consumului de utilități. Soluția implementată acoperă întregul traseu al datelor, de la generare până la vizualizare și predicție, și oferă o bază tehnică extensibilă pentru dezvoltări ulterioare în domeniul monitorizării inteligente a utilităților.
